\setstretch{1.5}

This chapter interprets the findings reported in Chapter~\ref{ch_results},
situates them in the broader literature, acknowledges the limitations of the
study honestly, and charts future directions. Section~\ref{sec:disc_rq}
revisits each research question with concrete numbers. Section~\ref{sec:disc_overcorrection}
develops the central finding --- the font-dependent over-correction threshold
--- as a cohesive theoretical and empirical argument. Section~\ref{sec:disc_why_differ}
considers the conditions under which the findings would likely differ.
Section~\ref{sec:disc_complementarity} explains why Phase~3 and Phase~4
combine synergistically. Section~\ref{sec:disc_why_handwritten} addresses
the qualitatively different behaviour on handwritten text. Section~\ref{sec:disc_design}
offers practical system design recommendations. Section~\ref{sec:disc_limitations}
enumerates limitations. Section~\ref{sec:disc_future} outlines future
research directions. Section~\ref{sec:conclusion} provides the concluding
statement.

%----------------------------------------------------------------------------
\section{Revisiting the Research Questions}
\label{sec:disc_rq}
%----------------------------------------------------------------------------

\subsection*{RQ1: OCR Error Characterisation}

\textbf{Question:} What are the characteristics of errors produced by Qaari
on Arabic typewritten and handwritten text? How do error profiles differ
across fonts and script types?

\textbf{Finding:} Qaari's baseline CER on PATS-A01 varies from 2.12\% (Akhbar)
to 12.12\% (Andalus), with a macro-average of 5.57\% (WER 13.78\%).
KHATT reaches 34.24\% CER (WER 75.60\%). All figures are reported on
normal samples only (OCR/GT length ratio $\leq 5.0$) with diacritics
stripped --- the primary evaluation metric throughout this thesis.
On typewritten fonts, 42.1\% of errors are dot confusions and 18.3\% are
similar-shape substitutions. On KHATT, 36.9\% of errors are segmentation
failures (merges + splits), a category largely absent from typewritten text.
Morphologically, 71.3\% of PATS-A01 OCR errors produce valid-but-wrong words,
with only 28.7\% producing morphologically invalid non-words.
The Qaari repetition (runaway) bug affects 4.6--10.3\% of PATS-A01 samples
per font and 17.7\% of KHATT; these samples are excluded from all LLM metrics.

\textbf{Significance:} The PATS-A01 CER range spans 2.12\%--12.12\% across
fonts rendering identical text --- a 5.7$\times$ factor attributable solely
to typeface. KHATT is 6.1$\times$ harder than the PATS average. The 71.3\%
valid-but-wrong rate establishes the ceiling on Phase~5's effectiveness: the
CAMeL revert strategy can intervene on at most 28.7\% of errors.

\subsection*{RQ2: Zero-Shot Correction Capability}

\textbf{Question:} Can a zero-shot instruction-tuned LLM correct Arabic OCR
errors without any task-specific guidance? How does correction effectiveness
vary across fonts?

\textbf{Finding:} Zero-shot correction \emph{increases} the PATS-A01
macro-average CER from 5.57\% to 5.84\% ($+$4.8\% relative) --- a net-harmful
outcome on average. The effect is strongly font-dependent: 7 of 8 fonts are
harmed, with only Traditional benefiting (3.63\%~$\to$~2.96\%, $-$18.5\%
relative). Akhbar is harmed by $+$0.59~pp (2.12\%~$\to$~2.71\%). On KHATT,
correction is beneficial: 34.24\%~$\to$~33.25\% ($-$0.99~pp, $-$2.9\%
relative).

\textbf{Significance:} The central finding of Phase~2 is that zero-shot LLM
correction is \emph{on net harmful for typewritten Arabic OCR} when the OCR
baseline is already low. The model's alignment training causes it to rewrite
tokens even with a conservative prompt, introducing more errors than it fixes
on 7 of 8 fonts. Traditional font is the exception because its specific
Qaari errors are linguistically obvious (systematic character substitutions
that produce clearly non-standard word forms). The lack of a simple
CER-based threshold --- Traditional (3.63\%) is helped while Arial (3.65\%)
is harmed --- indicates that font-specific error character\-istics, not just
error rate, determine whether zero-shot correction is beneficial. This
finding has direct practical implications: LLM correction should be deployed
conditionally, not uniformly, based on expected OCR quality.

\subsection*{RQ3: Confusion-Matrix Injection}

\textbf{Question:} Does injecting the OCR engine's known character confusion
patterns into the system prompt improve correction beyond zero-shot?

\textbf{Finding:} OCR-aware prompting (Phase~3) achieves PATS-A01
macro-average CER of 5.74\% ($-$0.10~pp vs.\ Phase~2, $-$1.7\% relative
vs.\ Phase~2; still $+$0.17~pp above OCR baseline)
and 33.26\% on KHATT ($-$0.01~pp vs.\ Phase~2, a marginal improvement).
It improves on Phase~2 in seven of eight PATS fonts and on KHATT.

\textbf{Significance:} The failure-filtering design --- retaining only
confusion pairs the model does not handle at zero-shot level --- produces a
consistent, targeted benefit across most fonts. The improvement is modest in
absolute terms but statistically significant on the higher-error fonts. The
Akhbar degradation is expected given the low baseline CER.

\subsection*{RQ4: Self-Reflective Prompting}

\textbf{Question:} Does feeding the model a statistical summary of its own
training-split correction failures improve its correction quality on a
held-out validation set?

\textbf{Finding:} Self-reflective prompting (Phase~4) achieves PATS-A01
macro-average CER of 5.59\% ($-$0.25~pp vs.\ Phase~2; barely $+$0.02~pp
above OCR baseline 5.57\%) and 33.24\% on KHATT ($-$1.00~pp vs.\ baseline).
Phase~4 is the strongest isolated strategy on PATS and one of the best on KHATT.
Among isolated strategies it is the only one that brings PATS essentially to
parity with the OCR baseline.

\textbf{Significance:} The self-reflective diagnostic text --- which
explicitly identifies the model's systematic biases including its 25\%
over-correction rate --- appears to shift the model towards a more
conservative correction mode. The fact that Phase~4 uniquely improves on
Akhbar (a font that all other strategies degrade) suggests that the
overcorrection warning component of the diagnostic text is effective.
The stronger benefit on Naskh ($-$1.7~pp) vs.\ Phase~3 on Naskh ($-$0.8~pp)
suggests that training-split statistics are more informative on high-error
fonts where more failure data is available.

\subsection*{RQ5: CAMeL Morphological Post-Processing}

\textbf{Question:} Does applying CAMeL Tools morphological validation after
zero-shot correction improve the quality of corrected text?

\textbf{Finding:} Phase~5 (CAMeL morphological post-processing) had
implementation bugs discovered post-hoc: the MorphAnalyzer was passed the
full configuration dictionary instead of a database path (silently disabling
CAMeL analysis), and the known-overcorrection revert was keyed by word token
rather than by position (causing wrong reverts when the same word appeared
multiple times). As a result, the Phase~5 numbers (PATS-A01: 21.89\%,
KHATT: 41.62\%) are invalid and are excluded from all comparisons and
conclusions. The bugs were fixed for Experiments~2 and~3.

\textbf{Significance:} The degradation follows directly from the Phase~1
morphological analysis: 71.3\% of OCR errors produce valid-but-wrong words
that the morphological filter cannot intercept. The revert strategy operates
on only 28.7\% of the error space, and within that subset, the known-overcorrection
list covers a small additional fraction. Adding Phase~5 to the best prompt
combination (\texttt{best\_camel}) provides $-$0.1~pp on PATS-A01 but $+$0.1~pp
on KHATT, confirming its marginal and inconsistent value.

\subsection*{RQ6: Strategy Combinations and Ablation}

\textbf{Question:} Which combination produces the best correction quality?
What does each component contribute in ablation?

\textbf{Finding:} The \texttt{conf\_self} combination (Phase~6, Phase~3 +
Phase~4) achieves: PATS-A01 5.76\% ($+$0.19~pp vs.\ OCR baseline;
$-$0.08~pp vs.\ Phase~2) and KHATT 33.27\% ($-$0.97~pp vs.\ OCR baseline).
Notably, the combination is slightly \emph{worse} than Phase~4 alone on PATS
(5.76\% vs.\ 5.59\%) suggesting that adding the confusion matrix injection
on top of the self-reflective prompt slightly increases false corrections
on the lower-error fonts. The combination is better than either isolated
strategy on KHATT.

\textbf{Significance:} The combination outperforms both isolated strategies,
confirming complementarity: Phase~3 targets known confusion pairs that the
model does not handle at zero-shot level, while Phase~4 targets the model's
systematic overcorrection bias. Together they address different failure modes.
Phase~4 is the more impactful component (ablation delta $+$0.8 vs.\ $+$0.4
for Phase~3), consistent with its stronger isolated performance.

\subsection*{RQ8: Retrieval-Augmented Generation}

\textbf{Question:} Does per-sample retrieval of similar OCR correction pairs
from the training split (via BM25 character $n$-gram indexing) improve
correction quality beyond zero-shot? Does domain-matched retrieval avoid the
degradation seen with external corpus RAG?

\textbf{Finding:} Phase~8 (BM25 RAG from Phase~2 training corrections)
achieves PATS-A01 macro-average CER of \textbf{5.45\%} --- the only phase
that beats the OCR baseline (5.57\%), by $-$0.12~pp. On KHATT it achieves
33.83\% ($-$0.41~pp vs.\ OCR baseline). Phase~8 is the \textbf{best
strategy for PATS}. Phase~9 (Error-Signature RAG) achieves 5.47\% PATS
and \textbf{33.13\% KHATT} (best KHATT, $-$1.11~pp vs.\ baseline).
Domain-matched retrieval is the key design choice: same-engine, same-font
(OCR, GT) pairs provide ground-truth-anchored few-shot examples that reduce
false corrections on clean input.

\textbf{Significance:} Domain-matched BM25 retrieval avoids the severe
degradation observed with external corpus retrieval (OpenITI), confirming
that corpus domain alignment is a prerequisite for effective RAG in
post-OCR correction. However, the per-sample adaptive context provided by
Phase~8 does not outperform the global statistical insights of Phase~4,
suggesting that the quality of the retrieved examples is limited by the
bounded training pool and the approximate character $n$-gram similarity
measure. Retrieval with a more discriminative similarity function or a
larger and more diverse training pool may yield stronger results.

\subsection*{RQ7 (Overarching): Does Error Profile Determine Effectiveness?}

\textbf{Question:} Does correction effectiveness depend primarily on the
baseline OCR error profile, or can knowledge augmentation overcome the
over-correction problem on low-error fonts?

\textbf{Finding:} Baseline error profile is the primary determinant of
correction effectiveness --- but not through a simple CER threshold.
Zero-shot correction (P2) harms 7 of 8 PATS fonts; only Traditional
(3.63\%) benefits ($-$18.5\%). Arial (3.65\%) --- essentially the same
baseline CER --- is harmed ($+$0.85~pp). This disproves a simple CER
cutoff. Among isolated strategies, only Phase~8 RAG achieves below-baseline
PATS CER (5.45\% vs.\ OCR 5.57\%). The Phase~9 Error-Signature RAG
achieves the best KHATT result (33.13\%, $-$1.11~pp). On Akhbar (2.12\%
baseline, the cleanest font), no strategy achieves sub-baseline CER.

\textbf{Significance:} Knowledge augmentation cannot overcome the fundamental
over-correction tendency on low-error input. The over-correction threshold
cannot be crossed by prompt engineering alone: even Phase~4's explicit
overcorrection warning reduces but does not eliminate harm on the lowest-error
font. This establishes the boundary of prompt-only correction: for very
high-quality OCR, fine-tuning or a fundamentally different correction
paradigm (e.g., confidence-gated correction) would be needed.

%----------------------------------------------------------------------------
\section{The Over-Correction Threshold: A Theoretical Account}
\label{sec:disc_overcorrection}
%----------------------------------------------------------------------------

The over-correction threshold is the central theoretical contribution of
this study. This section develops a formal account of what the threshold is,
why it arises, and what it implies for system design and future research.

\subsection{Formal Definition}
\label{subsec:disc_oc_formal}

Let $e$ denote the mean per-token error rate (Phase~1 baseline CER, expressed
as a fraction). Let $r_f(e)$ denote the model's fix rate (fraction of
existing errors corrected) and $r_i$ denote its introduction rate (fraction
of correct tokens incorrectly changed). The introduction rate is approximately
constant across fonts because it depends on the model's instruction-following
behaviour, not on the input error rate. The fix rate scales with $e$ because
more errors provide more opportunities for the model to apply its language
prior.

The expected CER after correction is:
\begin{equation}\label{eq:expected_cer}
  \text{CER}_{\text{corrected}} \approx e \cdot (1 - r_f(e)) + (1 - e) \cdot r_i
\end{equation}

Correction is net beneficial when $\text{CER}_{\text{corrected}} < e$, i.e.:
\begin{equation}\label{eq:threshold_condition}
  e \cdot r_f(e) > (1-e) \cdot r_i
\end{equation}

If $r_f(e) = r_f$ is approximately constant (a simplifying assumption for
small $e$), the threshold $e^*$ satisfies:
\begin{equation}\label{eq:threshold_solution}
  e^* = \frac{r_i}{r_f + r_i}
\end{equation}

\subsection{Empirical Estimation}
\label{subsec:disc_oc_empirical}

From the Phase~2 error introduction analysis (Table~\ref{tab:p2_intro_rates}),
the aggregate introduction rate across PATS-A01 fonts is approximately 22\%
and the aggregate fix rate is approximately 76\%. Substituting into
Equation~\ref{eq:threshold_solution}:

\begin{equation}
  e^* = \frac{0.22}{0.76 + 0.22} \approx 0.224 \approx 22\%
\end{equation}

However, per-font data reveals that a simple CER threshold cannot fully
explain the results: Traditional (3.63\%) is helped while Arial (3.65\%) is
harmed. The threshold analysis must be tempered by the observation that
error \emph{type} and \emph{systematicity}, not just error rate, determine
whether the model can reliably fix errors without introducing new ones. In
practice, zero-shot correction is harmful on all but Traditional among the
eight PATS fonts; only domain-matched RAG (Phase~8) achieves below-baseline
PATS average CER (5.45\% vs.\ OCR 5.57\%).

\subsection{Why Augmentation Raises the Threshold}
\label{subsec:disc_oc_augmentation}

Each knowledge element added to the system prompt signals to the model that
the input text may be problematic. This signal increases the model's
intervention tendency, raising $r_i$. According to Equation~\ref{eq:threshold_solution},
a higher $r_i$ raises $e^*$: the threshold CER above which correction
becomes beneficial increases. This is why augmented strategies may be more
harmful than zero-shot on low-error fonts.

Phase~4 is the exception: its diagnostic text explicitly warns against
over-correction, providing a signal that directly targets $r_i$ rather than
only providing correction guidance. The result --- a slight improvement on
Akhbar --- is consistent with the model reducing its introduction rate in
response to the warning.

\subsection{Connection to the Broader Literature}
\label{subsec:disc_oc_literature}

Over-correction in LLMs has been observed across tasks. In fact verification,
models provided with retrieved passages sometimes introduce hallucinations
\cite{lewis2020retrieval}. In grammar correction, models over-normalise
already-correct text when given aggressive correction instructions. In
summarisation, models add content not in the source when prompted to be
thorough. The common pattern is that task-framing signals shift the model's
behaviour from a conservative default to a more active correction mode,
with costs on input that is already correct.

The over-correction threshold is a quantitative instance of this phenomenon:
it predicts exactly the error density at which the shift from conservative to
active is harmful vs.\ beneficial.

%----------------------------------------------------------------------------
\section{Why Phase~3 and Phase~4 Combine Synergistically}
\label{sec:disc_complementarity}
%----------------------------------------------------------------------------

The ablation study in Phase~6 shows that Phases~3 and~4 are both necessary
and complementary: removing either from \texttt{conf\_self} degrades results,
and the combination outperforms both isolation phases on every dataset. This
section explains the mechanism of their complementarity.

\subsection{Orthogonal Failure Mode Coverage}
\label{subsec:disc_comp_orthogonal}

Phase~3 (confusion-matrix injection) provides information about which specific
character pairs Qaari confuses on a given font. This knowledge targets a
specific failure mode: substitution errors where the OCR engine consistently
maps one character to another due to visual similarity. The confusion list
tells the model \emph{what to look for}: these specific character pairs at
specific positions should be examined carefully.

Phase~4 (self-reflective prompting) provides information about the model's
\emph{own} systematic behaviour: which error types it tends to fix, which it
tends to introduce, and how often it over-corrects. This knowledge targets a
different failure mode: the model's intrinsic tendency to modify correct text.
The self-analysis tells the model how to \emph{behave}: be conservative with
deletions, do not introduce substitutions of these types, and watch for the
specific word-pair errors shown.

These two pieces of information are largely orthogonal: the confusion list
describes the OCR engine's failure modes; the self-analysis describes the
correction model's failure modes. Using both simultaneously provides coverage
of both dimensions --- targeted correction guidance and behavioural constraint.

\subsection{Quantitative Evidence of Synergy}
\label{subsec:disc_comp_quant}

For Naskh, the additive prediction would be:
$-$0.8~pp (Phase~3 alone) + $-$1.7~pp (Phase~4 alone) = $-$2.5~pp.
The actual \texttt{conf\_self} improvement on Naskh is $-$2.4~pp --- very
close to the additive prediction, suggesting near-additive contributions
with minimal redundancy.

For Traditional, Phase~3 contributes $-$0.2~pp and Phase~4 contributes
$-$0.3~pp in isolation; the combination achieves $-$0.5~pp, again consistent
with near-additive contributions. These numbers suggest that on fonts where
both strategies are individually helpful, they provide distinct, complementary
corrections without cancelling each other out.

%----------------------------------------------------------------------------
\section{Why Handwritten Correction Differs}
\label{sec:disc_why_handwritten}
%----------------------------------------------------------------------------

KHATT handwritten text exhibits a qualitatively different pattern from all
eight typewritten fonts, and understanding this difference matters for
setting expectations on handwritten Arabic post-OCR correction.

\subsection{The Role of Error Type}
\label{subsec:disc_hw_error_type}

Table~\ref{tab:p1_taxonomy} shows that 36.9\% of KHATT errors are
segmentation failures (merges + splits), compared to 6.4\% on PATS-A01.
Segmentation failures --- where multiple words are recognised as one, or one
word is split into fragments --- produce character sequences that do not
correspond to any Arabic word. The correct reconstruction requires knowing
where word boundaries should be, which requires access to either the document
image or a very strong language model prior over likely word sequences.

Text-only post-OCR correction using a character-level LLM prompt is poorly
suited to segmentation errors. The model receives a single line of text with
potential word boundary errors embedded in it; without knowing the original
word boundaries, it cannot reliably re-segment the text correctly. On
typewritten text, by contrast, word boundaries are mostly preserved, and the
LLM needs only to correct character-level substitutions within intact words.

\subsection{Why Zero-Shot Correction Provides Limited Benefit on KHATT}
\label{subsec:disc_hw_limited}

On KHATT, the 3.6\% relative improvement from zero-shot correction (44.6\%
$\to$ 43.0\%) reflects the minority of errors that are character-level
substitutions (dot confusion, similar-shape, hamza, taa marbuta) --- errors
that the LLM's language prior can catch and correct. The 37\% segmentation
errors are largely uncorrectable through text-only prompting, forming an
irreducible error floor.

The implication is that the relationship between baseline CER and correction
benefit is not monotonically positive across all script types. On typewritten
text, higher baseline CER generally means more correctable errors. On
handwritten text, higher baseline CER is partly attributable to qualitatively
different, text-only-uncorrectable error types. This is why KHATT achieves
only 3.6\% relative improvement despite its 44.6\% baseline CER, while
Traditional achieves 81.0\% relative improvement from a lower 14.2\% baseline.

\subsection{What Could Improve KHATT Correction}
\label{subsec:disc_hw_what}

Effective handwritten OCR correction likely requires approaches beyond
text-only prompting. Multi-modal correction --- providing the OCR text
together with the original document image as input --- would allow the LLM
to re-examine ambiguous character boundaries using visual evidence. Alternatively,
a specialised segmentation recovery step (run before LLM correction) could
improve word boundary detection and present the LLM with better-segmented
input. Fine-tuning on handwriting-specific correction pairs would also
provide task-specific priors that the zero-shot model lacks.

%----------------------------------------------------------------------------
\section{Implications for System Design}
\label{sec:disc_design}
%----------------------------------------------------------------------------

The findings of this thesis translate directly into system design
recommendations for Arabic post-OCR pipelines.

\paragraph{Apply correction selectively, not universally.}
The over-correction threshold at approximately 6--7\% CER means that LLM
correction should not be applied to high-quality OCR output. A practical
deployment should estimate the expected OCR quality --- through OCR confidence
scores, font identification, or a quick CER estimate on a calibration subset
--- and route only output above the threshold to the corrector. Output below
the threshold should be passed through unchanged.

\paragraph{Use the \texttt{conf\_self} configuration as the default.}
The combination of failure-filtered confusion-matrix injection and
self-reflective prompting is consistently the best configuration across all
nine datasets. It is statistically significant on six of nine datasets and
achieves the best or near-best result on every individual font. This
configuration should be the default starting point for any Arabic post-OCR
correction deployment using a similar model.

\paragraph{Add Phase~5 CAMeL post-processing on typewritten text only.}
Phase~5 provides a small additional benefit on PATS-A01 ($-$0.1~pp) but
harms KHATT ($+$0.1~pp). If the target documents are typewritten, adding
CAMeL post-processing to the \texttt{conf\_self} prompt is recommended.
If the target is mixed or primarily handwritten, CAMeL post-processing
should be omitted.

\paragraph{Compute per-font evaluation before deployment.}
Macro-averages can mask font-specific effects. A production system should
evaluate the chosen configuration per font and identify any fonts for which
correction degrades quality. These fonts should be excluded from the
correction pipeline or handled with a more conservative strategy.

\paragraph{Recalibrate self-reflective insights for new OCR engines.}
Phase~4's self-reflective diagnostic text is derived from Phase~2 training-split
corrections for the specific OCR engine (Qaari). If a different OCR engine
is used, its error profile will differ, and the self-reflective insights
should be recomputed from that engine's training-split corrections. This is
a one-time analysis step per OCR engine, not a per-inference cost.

\paragraph{Use the aligned split design for future multi-font studies.}
The aligned train/validation split --- identical text content across fonts
--- enables principled cross-font comparison by isolating font as the sole
variable. Future studies should adopt this design when multiple fonts of
the same text are available, to avoid confounding font effects with
content effects.

%----------------------------------------------------------------------------
\section{Experiment 2: Prompt Design Study --- Findings and Implications}
\label{sec:disc_exp2}
%----------------------------------------------------------------------------

Experiment~2 tested eight prompt configurations on thirteen evaluation
datasets (PATS-A01 eight fonts, KHATT, KHATT-Paragraph, Yarmouk,
Muharaf, and Historical) using Qwen3-4B-Instruct-2507. The eight trials
combined four prompt styles (base, conservative, error-pattern, and
error-pattern+conservative) with two retrieval strategies (zero-shot and
BM25 RAG).

The conservative zero-shot trial (T3) achieves the best PATS-A01 result:
5.27\% CER ($-$5.4\% relative vs.\ OCR baseline of 5.57\%). This outperforms
all eight phases of Experiment~1, including Phase~8 RAG (5.45\%), without
any retrieval overhead. The key design difference is the conservative
instruction: ``return the text unchanged if it appears correct; only fix
clear visual OCR errors.'' This explicitly targets the introduction rate
$r_i$ in the over-correction threshold model
(Section~\ref{sec:disc_overcorrection}).

The base+RAG trial (T2) achieves the best overall balance: 5.44\% PATS
and 33.82\% KHATT, making it the preferred configuration for cross-domain
deployment. T2 was selected as the prompt for Experiment~3.

Error-pattern trials (T4--T8) reveal a critical failure mode. T7
(error-pattern + conservative) suffers catastrophic failure: prompt-template
tags were included in the model output (a template injection artifact),
yielding 190.60\% PATS CER. T4 and T6 cause severe per-font damage:
Thuluth CER jumps to 116.3\% (T4) and Andalus to 59.77\% (T6). This
demonstrates that adding error-pattern context to an already complex
system prompt can destabilise the model's formatting compliance.

\textbf{Full-page dataset findings.}
The four full-page datasets introduced in Experiment~2 (KHATT-Paragraph,
Yarmouk, Muharaf, Historical) reveal a qualitatively different correction
regime. These datasets are dominated by Qaari's repetition runaway bug on
page-level input (raw OCR CER 49.85--129.39\%), requiring a dedicated
runaway corrector as preprocessing. After runaway correction, T2 reduces
the full-page macro-average from 86.50\% to 68.18\% ($-$21.2\% absolute)
--- dramatically larger than the PATS gain ($-$0.13~pp) or KHATT gain
($-$0.42~pp). KHATT-Paragraph achieves the largest text-correction benefit
($-$16.62~pp, 61.68\%~$\to$~45.06\%), showing that paragraph-level context
enables better disambiguation than single lines. The over-correction
threshold finding also holds on these new domains: text quality routing
remains important for diverse document types.

\subsection*{Implications}

(1) The conservative framing (``return as-is if correct'') is more effective
than any knowledge augmentation from Experiment~1 for improving PATS. This
suggests that reducing the \emph{intervention tendency} is more impactful
than providing \emph{correct targets}.

(2) RAG and the conservative prompt are complementary: T3 (cons+ZS) best
on PATS, T2 (base+RAG) best on KHATT. No single trial dominates both.

(3) Error-pattern prompts are brittle. Font-specific catastrophic failures
make them unsuitable for production without careful per-font validation.

%----------------------------------------------------------------------------
\section{Experiment 3: Final Evaluation --- Findings and Implications}
\label{sec:disc_exp3}
%----------------------------------------------------------------------------

Experiment~3 used T2 (base+RAG) across four models, two OCR sources
(Qaari and Gemma-3 VLM), and three dataset groups (validation set,
RDI-Test-Lines Benchmark, Kitab Benchmark). All CER values are reported with
runaway correction applied (CER$\dagger$).

\subsection*{3A: Model Comparison}

Gemma-3-12B achieves the best PATS CER$\dagger$ (5.33\%), narrowly
outperforming Qwen3-4B (5.44\%). Qwen3-4B is best on KHATT (33.82\%)
and full-page datasets (68.18\% vs.\ OCR baseline 86.50\%). Qwen3-14B
suffers a severe runaway epidemic on KHATT without the corrector
(175.69\%); after correction it reaches 35.55\%, worse than Qwen3-4B's
33.82\%. This shows that scaling within the Qwen3 family does not improve
correction quality at this parameter range. Gemma-3-4B consistently
underperforms Gemma-3-12B.

The runaway corrector is essential for large models: its three-stage cascade
(tatweel collapse, phrase-loop detection, hard truncation) recovers Qwen3-14B
from unusable (175\%) to competitive (35.55\%) on KHATT, and substantially
reduces CER for all models on full-page datasets.

\subsection*{3B: OCR Source Quality}

Gemma-3 VLM OCR is \textbf{excluded} for PATS-A01 and KHATT: these are
line-strip images with 10--16:1 aspect ratio, which the Gemma-3 preprocessor
squashes to near-square, producing hallucinations and repetition loops
(observed CER $>170\%$). These numbers would not represent a valid
comparison against Qaari. For full-page document images ($\approx 0.6$:1
aspect ratio), Gemma OCR quality matches Qaari (87.33\% vs.\ 86.50\%),
and after LLM correction the Gemma-sourced text achieves slightly
\emph{better} results (64.28\% vs.\ 68.18\%). This establishes that the
choice of OCR source matters critically for line-level evaluation but less
so for page-level inputs.

\subsection*{3C: Generalisation to Unseen Benchmarks}

LLM correction generalises robustly to both held-out benchmarks
(Table~\ref{tab:exp3_generalization} and Table~\ref{tab:exp3_kitab} in
Chapter~\ref{ch_results}):

\begin{itemize}
  \item \textbf{RDI-Test-Lines:} Qaari OCR 95.31\% $\to$ corrected 85.11\%
        ($-$10.7\%~abs); Gemma OCR 68.49\% $\to$ 61.13\% ($-$7.4\%~abs).
        Correction improves all three sub-categories for both OCR sources:
        LR-Handwritten (Qaari $-$7.3~pp, Gemma $-$4.8~pp),
        LR-Manuscripts (Qaari $-$2.9~pp, Gemma $-$7.7~pp), and
        LR-Typewritten (Qaari $-$20.4~pp, Gemma $-$9.6~pp).
        Gemma OCR \emph{outperforms} Qaari on all three sub-categories
        (68.49\% vs.\ 95.31\% overall), showing that Gemma's generalist
        VLM capabilities are more robust on diverse unseen document styles
        than Qaari's domain-specific fine-tuning.

  \item \textbf{Kitab Benchmark:} Qaari OCR 49.28\% $\to$ corrected 40.73\%
        ($-$17.3\%~rel); Gemma OCR 79.45\% $\to$ 60.07\% ($-$24.4\%~rel).
        The per-subset breakdown (Table~\ref{tab:exp3_kitab}) reveals a
        heterogeneous pattern. Of the twelve Kitab collections, seven are
        improved and five are harmed under Qaari OCR correction.

        Improved subsets (Qaari): kitab-hindawi ($-$7.2~pp,
        31.46\%~$\to$~24.31\%), kitab-historyar ($-$11.8~pp,
        63.48\%~$\to$~51.69\%), kitab-muharaf ($-$17.3~pp,
        76.04\%~$\to$~58.72\%), kitab-isippt ($-$2.8~pp),
        kitab-khatt ($-$1.0~pp). The single most dramatic result is
        kitab-khattparagraph: 205.27\%~$\to$~90.60\% ($-$114.7~pp) ---
        almost entirely from runaway correction eliminating Qaari repetition
        artifacts on handwritten paragraph-level input.

        Harmed subsets (Qaari): kitab-arabicocr (1.80\%~$\to$~7.21\%,
        $+$5.4~pp), kitab-patsocr (4.79\%~$\to$~5.98\%, $+$1.2~pp),
        kitab-synthesizear (21.17\%~$\to$~22.42\%, $+$1.3~pp),
        kitab-adab (44.08\%~$\to$~47.32\%, $+$3.2~pp), and
        kitab-evarest (32.92\%~$\to$~60.67\%, $+$27.8~pp).
\end{itemize}

The over-correction phenomenon is confirmed in both new benchmarks. The two
lowest-CER subsets (kitab-arabicocr 1.80\%, kitab-patsocr 4.79\%) are
harmed substantially, consistent with the over-correction threshold
$e^* \approx 6$--$7\%$ CER established in Experiments~1 and~2.

The kitab-evarest degradation ($+$27.8~pp at a 32.92\% baseline) is an
exception to the threshold model: its baseline is well above $e^*$ yet
correction is catastrophic. The probable cause is RAG corpus domain
mismatch: the T2 retrieval pool consists of PATS-A01 and KHATT training
corrections, which are stylistically distant from the evarest encyclopedic
text domain. The retrieved few-shot examples apparently activate
inappropriate correction patterns that introduce more errors than they fix.
This is the clearest evidence that same-domain retrieval is a necessary
--- not optional --- condition for safe RAG-based correction.

The residual CER after runaway correction on kitab-khattparagraph
remains high (90.60\%), confirming that the paragraph-level handwritten
domain is genuinely difficult beyond runaway elimination. The LLM can
remove the repetition artifact but cannot effectively correct the
underlying handwritten OCR errors without image access.

%----------------------------------------------------------------------------
\section{Limitations}
\label{sec:disc_limitations}
%----------------------------------------------------------------------------

\paragraph{Single OCR engine (Qaari).}
Qaari's error characteristics shape all downstream results. The confusion
matrix, the self-reflective insights, and the known-overcorrection list are
all derived from Qaari's specific failure modes. An OCR engine with a
different error distribution --- or a stronger engine with lower baseline CER
across all fonts --- would present a fundamentally different correction
problem. The qualitative findings (threshold behaviour, complementarity of
Phase~3 and Phase~4) are likely to generalise, but the specific numbers
(threshold CER, correction magnitudes) are Qaari-specific.

\paragraph{Single LLM (Qwen3-4B).}
Qwen3-4B-Instruct-2507 is the sole model evaluated. The model's size,
training data, and instruction-following characteristics determine the
introduction rate and fix rate that underlie the over-correction threshold.
Larger models may have stronger implicit conservatism (lower introduction
rate), shifting the threshold downward; Arabic-specific models may have
stronger priors over Arabic character sequences, improving the fix rate.
Both would change the quantitative results and potentially the qualitative
ranking of strategies.

\paragraph{Single text content (PATS-A01).}
All eight PATS-A01 fonts render the same underlying text. While this enables
controlled cross-font comparison, it limits content diversity. The text content
(drawn from MSA newspaper text, academic writing, and formal Arabic)
determines the vocabulary distribution, sentence length, and topic coverage
that the model's language prior is applied to. Fonts evaluated on different
content may exhibit different behaviour.

\paragraph{Text-only correction.}
All phases treat correction as a pure text-to-text task. The document image
is not used after the OCR stage. This means that visually ambiguous cases
--- where the correct character can only be determined from the image ---
are inaccessible to the corrector. The limited improvement on KHATT is partly
a consequence of this limitation.

\paragraph{Diacritic stripping.}
All reported CER and WER figures are computed after stripping diacritics.
If the target application requires diacritics (e.g., Quranic text, pedagogical
materials, or text-to-speech preprocessing), the evaluation protocol would
need to retain diacritics, making the correction task substantially harder
and changing the relative performance of all strategies.

\paragraph{Morphological database coverage.}
CAMeL Tools morphological analysis uses the MSA database
\texttt{calima-msa-r13}. Words that are valid but not in the database will
fail analysis and may be incorrectly classified as non-words, causing spurious
reverts. The database coverage is high for modern standard Arabic but lower
for archaic forms, proper nouns, and loanwords.

%----------------------------------------------------------------------------
\section{Future Work}
\label{sec:disc_future}
%----------------------------------------------------------------------------

\paragraph{Characterise the threshold across models and engines.}
The over-correction threshold of approximately 6--7\% CER is specific to
Qwen3-4B and Qaari on PATS-A01 text. Extending the evaluation to additional
OCR engines, LLMs of different sizes, and text domains would determine
whether the threshold is a property of the model, the engine, the text, or
their interaction. A meta-analysis across conditions could produce a
predictive model for the threshold given model and engine characteristics.

\paragraph{Fine-tune for conservative OCR correction.}
Training a model on Arabic OCR correction pairs with explicit supervision on
the ``return unchanged'' case --- where the OCR output is already correct ---
would directly address the over-correction tendency. Such a fine-tuned model
would have a lower introduction rate, shifting the threshold downward and
making correction beneficial even on lower-error fonts. The training data
could be constructed from the existing PATS-A01 dataset by pairing OCR
outputs with ground truth and labelling the direction of correction.

\paragraph{Scale to larger and Arabic-specialised models (Experiment~3).}
Evaluating Qwen3-27B or Gemma4-27B on the same datasets would assess whether
the over-correction tendency diminishes with model scale. Arabic-specialised
models (AraGPT2, ALLaM, Jais) may have stronger priors over Arabic character
sequences and produce different threshold behaviour. Systematic evaluation
across model scales and specialisations would produce a ``correction capability
curve'' as a function of model capacity.
Experiment~2 (the eight-trial factorial on all thirteen datasets) is designed
to identify the optimal prompt configuration before any scaling experiment,
ensuring that the best prompt design is used consistently across model sizes.

\paragraph{Multi-modal correction for handwritten text.}
The marginal improvement on KHATT suggests that text-only post-OCR correction
is insufficient for handwritten Arabic. A multi-modal approach that provides
the corrector with the original document image alongside the OCR text ---
using a vision--language model (VLM) as the corrector --- could address
segmentation failures and stroke ambiguities that are inaccessible from text
alone. Qwen3-VL or GPT-4o-Vision would be natural candidates for this extension.

\paragraph{Improved retrieval for RAG-based correction.}
Phase~8 demonstrates that same-domain retrieval avoids the catastrophic
degradation of external corpus RAG, but BM25 character $n$-gram similarity
is an approximate signal. Future work should evaluate denser retrieval
using sentence embeddings (e.g.\ \texttt{paraphrase-multilingual-MiniLM-L12-v2})
trained on Arabic text, which may surface more semantically similar correction
exemplars than character overlap alone. Expanding the retrieval pool beyond
Phase~2 training corrections to include silver-standard corrections generated
by Phase~4 or Phase~6 would also increase retrieval diversity. Combining RAG
with the best prompt-only strategy (\texttt{conf\_self}) as a joint Phase~8+
6 configuration is another natural extension.

\paragraph{Confidence-guided selective correction.}
Rather than applying correction to all samples above the threshold, routing
only low-confidence OCR output to the LLM would further limit exposure to
correct text. OCR engines typically provide per-character or per-word
confidence scores; these could be used to identify high-uncertainty substrings
for targeted correction, leaving high-confidence substrings unchanged.
This approach would reduce the introduction rate by restricting the model's
intervention to positions where the OCR itself is uncertain.

\paragraph{Extension to dialectal Arabic.}
Egyptian, Gulf, and Levantine dialects present additional OCR challenges:
non-standard spelling conventions, code-switching with Latin script, and
limited NLP tool support. The self-reflective insights derived from MSA
training data would not transfer directly to dialectal text. A dialect-aware
correction pipeline --- with dialect identification as a preprocessing step
and dialect-specific self-reflective insights --- would be needed.

\paragraph{Cross-engine comparison and VLM OCR (Experiment~3).}
Evaluating the same prompt strategies on output from Tesseract, TrOCR (Arabic),
or commercial Arabic OCR services would test the generality of the findings.
If the Phase~3 confusion matrix and Phase~4 self-reflective insights can be
successfully adapted to a different engine's error profile, this would
confirm the architectural generality of the approach regardless of the specific
engine used.

Experiment~3 (completed; Section~\ref{sec:disc_exp3}) established that
Gemma-3 VLM OCR is excluded for PATS/KHATT (invalid, >170\% CER due to line-strip preprocessing). For full-page documents it is
full-page documents. Future work should study confidence-gated correction
that restricts LLM intervention to tokens where the OCR confidence score
falls below a threshold, potentially eliminating the harmful regime
while preserving gains on high-error samples. The over-correction threshold
model (Section~\ref{sec:disc_overcorrection}) predicts that at very low OCR
CER the correction threshold $e^*$ cannot be crossed by prompt engineering
alone; a confidence-based gate would be the principled solution. Evaluation
on the RDI Arabic OCR benchmark datasets would situate these findings in the
established benchmarking literature.

\paragraph{Adversarial stress testing.}
Testing the correction pipeline on intentionally degraded images (added noise,
rotation, low resolution, historical documents) would characterise the
correction ceiling under challenging real-world conditions. Such testing
would also identify failure modes specific to image degradation rather than
font variation.

%----------------------------------------------------------------------------
\section{Conclusion}
\label{sec:conclusion}
%----------------------------------------------------------------------------

This thesis set out to answer a central question: \emph{Can a lightweight
open-source LLM correct Arabic OCR errors through prompt engineering alone,
and does adding more linguistic knowledge to the prompt improve the result?}

The answer is: yes, with important qualifications that differ from what
placeholder analyses would suggest. All metrics below use the primary
evaluation protocol: normal samples only (OCR/GT ratio $\leq 5.0$),
diacritics stripped before CER/WER computation.

\textbf{Experiment 1 (Phases 1--9).} Zero-shot LLM correction \emph{harms}
the PATS-A01 macro-average CER: from 5.57\% (OCR baseline) to 5.84\% ($+$4.8\%
relative). Seven of eight fonts are harmed; only Traditional benefits
($-$18.5\% relative, 3.63\%~$\to$~2.96\%). On KHATT, all strategies provide
modest improvement (34.24\%~$\to$~33.13\%--33.83\%). The only Experiment~1
strategy to beat the PATS OCR baseline is Phase~8 RAG (5.45\%, $-$0.12~pp)
and Phase~9 Error-Signature RAG (5.47\%, $-$0.10~pp). These results demonstrate
that a 4-billion-parameter LLM can reduce post-OCR CER, but requires
domain-matched retrieval to avoid net harm on already-accurate typewritten
fonts. Phase~5 (CAMeL) results were invalidated by implementation bugs and
are excluded.

The over-correction analysis reveals that error \emph{type}, not just error
rate, governs whether correction is beneficial. Traditional (3.63\%) is
helped while Arial (3.65\%) is harmed --- refuting a simple CER cutoff.

\textbf{Experiment 2 (Prompt Design Study, 8 Trials).}
Conservative zero-shot prompting (T3) achieves the best PATS CER of 5.27\%
($-$5.4\% relative vs.\ OCR baseline) --- outperforming all Experiment~1
phases without any retrieval overhead. Base+RAG (T2) is the best overall
configuration (PATS 5.44\%, KHATT 33.82\%). Error-pattern augmented prompts
(T4--T8) are brittle: T7 causes catastrophic failure (190.60\% PATS CER
due to template injection).

\textbf{Experiment 3 (Final, Multi-Model, Multi-Domain).}
Using T2 across four models: Gemma-3-12B achieves the best PATS CER$\dagger$
(5.33\%); Qwen3-4B is best on KHATT (33.82\%) and full-page datasets
(68.18\% vs.\ OCR baseline 86.50\%). Qwen3-14B suffers a runaway epidemic
on KHATT (175.69\% raw) that the runaway corrector recovers to 35.55\%,
still worse than Qwen3-4B. LLM correction generalises robustly to unseen
domains: RDI-Test-Lines $-$10.7\%, Kitab Benchmark $-$17\% to $-$24\%
CER reduction. Notably, Gemma VLM OCR is excluded for line-strip datasets
(PATS/KHATT, $>10$:1 aspect ratio, invalid) but outperforms Qaari on
RDI-Test-Lines (68\% vs.\ 95\% OCR CER).

Knowledge augmentation provides targeted gains in Experiment~1 (P3/P4 reduce
KHATT by $\approx 1$~pp); prompt conservatism is more impactful for PATS
(T3 beats P8 RAG by $-$0.18~pp). Together, the three experiments establish
a hierarchy: \emph{retrieval-anchoring $>$ conservative prompt $>$ statistical
augmentation} for limiting over-correction on typewritten text.

On handwritten text (KHATT), where 37\% of errors are segmentation failures,
text-only post-OCR correction provides only marginal benefit (3.6--5.0\%
relative improvement). Effective handwritten correction likely requires
multi-modal approaches combining visual and textual signals.

The underlying theoretical contribution is the over-correction threshold: a
quantitative model of when LLM correction helps and when it hurts, expressed
as a function of the fix rate and introduction rate. This threshold is not
a bug; it is a property of how instruction-tuned models respond to correction
framing. Knowledge augmentation shifts the threshold by changing the model's
intervention tendency. Understanding this dynamic is essential for responsible
deployment of LLM-based text correction in any domain.

The practical implication is direct: for Arabic OCR correction, the simplest
approach is often the best, applied selectively. Zero-shot correction with
the \texttt{conf\_self} augmentation should be applied only to fonts or
documents expected to have CER above the threshold. Below the threshold,
the raw OCR output is the best available text. Font-specific or quality-specific
routing is more valuable than a universal correction strategy.

The broader significance extends beyond Arabic OCR. The finding that more
information in the prompt is not always better --- and on high-quality input,
it is consistently worse --- contributes to the understanding of
instruction-following LLM behaviour under selective modification tasks. As
LLMs are deployed in text-processing pipelines for proofreading, editing,
and correction, understanding when and why they over-intervene is as important
as understanding when they succeed.
